\scanpage{45}
\mspage{4}
% L1–6
et les relations symétriques
\[
\left\{\begin{aligned}
\unsure{\zeta|\xi}_i&=2\eta'-\zeta'+\zeta'_i=2\eta'-\sum_{\alpha\ne i}\zeta'_i\\
&\deleted{\unsure{\eta=5\eta'-2\zeta'|\eta=5\eta'-2\eta'}}\\
\zeta&=12\eta'-5\zeta'
\end{aligned}\right.\tag*{(d')}
\]
d'où
\[
\zeta+\zeta'=4\lambda=4\lambda'\quad\text{avec}\quad
\lambda=3\eta-\zeta,\ \lambda'=3\eta'-\zeta'\tag*{(e)}
\]
\deleted{Donc}, On aura \deleted{\unsure{par} $4\unsure{(\lambda|(\mathrm e)}$} $4\lambda=4\lambda'$ i.e. $4(\lambda'-\lambda)=0$, mais on suppose qu'on a $=$
\[
3\eta-\zeta=3\eta'-\zeta'\overset{\mathrm{dfn}}=\lambda\tag*{\textcircled{f}}
\]
\margin{\uline{NB} conséquence de ce qui précède si 2 inv. dans $k$ (ou ds $M$)}
% L7–12
et on trouve \inserted{$3\eta'=3\eta-\zeta+\zeta'=15\eta-6\zeta$ compte tenu de \textcircled{b} donnant $2\eta'\deleted{\unsure{+\zeta}}=\zeta+\zeta'-2\eta=10\eta-4\zeta$}
\[
\left\{\begin{aligned}\eta+\eta'&=2\lambda\\\zeta+\zeta'&=4\lambda\end{aligned}\right.\tag*{(f)}
\]
\[
\left\{\begin{aligned}\eta'&=5\eta-2\zeta\\\eta&=5\eta'-2\zeta'\end{aligned}\right.\tag*{(g,g')}
\]
On voit que pour que $(\eta,(\zeta_i))$ soit une famille gén. (basique) \unsure{de|dans} $M$, il f et s qu'il en soit de $=$ de $(\eta',(\zeta'_i))$. \deleted{\unsure{Posant|Posons}} \inserted{On a}
\[
\eta-\zeta_i-\zeta_j=\eta'-\zeta'_i-\zeta'_j\quad(\overset{\mathrm{dfn}}=\zeta_{ij})\tag*{(\deleted{\unsure{f}}h)}
\]
% L13–20
Le module engendré par \deleted{\unsure{l'ens des}} $(\eta,(\zeta_i))$ ou par $(\eta',(\zeta'_i))$, ou par la \deleted{\unsure{2}} famille de 27 él. $(\zeta_i)$ $(\zeta'_i)$ $(\zeta_{ij})$ [ou \unsure{seult|même} $((\zeta_i),\zeta_{12})$ disons] sont les $=\text{s}$\ldots

Supposons qu'on \deleted{existe} \inserted{donne} sur $M$ une forme bil. \inserted{sym.}, notée $x\cdot y$. \deleted{telle que} \inserted{Alors} les relations
\[
\eta^2=1,\quad\zeta_i^2=-1,\quad\eta\zeta_i=0,\quad\zeta_i\zeta_j=0\quad\text{si }i\ne j\tag*{(i)}
\]
sont équiv. aux relations sym.
\[
\eta'{}^2=1,\ \zeta'_i{}^2=-1,\quad\eta'\zeta'_i=0,\quad\zeta'_i\zeta'_j=0\quad\text{si }i\ne j\tag*{(\deleted{\unsure{g}}i')}
\]
et elles impliquent que $(\eta,(\zeta_i))$ font une famille libre (ou $(\eta',(\zeta'_i))$ aussi) -- donc
% archive sheet number: 44
